\documentclass[11pt, letterpaper]{article}

% ── Packages ──────────────────────────────────────────────────────────────────
\usepackage[margin=1in]{geometry}
\usepackage{amsmath, amssymb}
\usepackage{hyperref}
\usepackage{xcolor}
\usepackage{booktabs}
\usepackage{array}
\usepackage{enumitem}
\usepackage{titlesec}
\usepackage{parskip}
\usepackage{fancyhdr}
\usepackage{siunitx}
\usepackage{longtable}
\usepackage[T1]{fontenc}

% ── Colour palette ─────────────────────────────────────────────────────────────
\definecolor{trinityblue}{RGB}{10, 60, 130}
\definecolor{darkgray}{RGB}{60, 60, 60}
\definecolor{lightgray}{RGB}{245, 245, 245}
\definecolor{accentblue}{RGB}{30, 100, 200}

% ── Hyperref setup ─────────────────────────────────────────────────────────────
\hypersetup{
  colorlinks=true,
  linkcolor=trinityblue,
  urlcolor=accentblue,
  citecolor=trinityblue
}

% ── Section title formatting ───────────────────────────────────────────────────
\titleformat{\section}
{\large\bfseries\color{trinityblue}}
{\thesection}
{0.6em}
{}
[\vspace{-0.5ex}\rule{\linewidth}{0.8pt}\vspace{0.5ex}]

\titleformat{\subsection}
{\normalsize\bfseries\color{darkgray}}
{\thesubsection}
{0.6em}
{}

\titlespacing{\section}{0pt}{1.2em}{0.5em}
\titlespacing{\subsection}{0pt}{0.8em}{0.3em}

% ── Header / Footer ────────────────────────────────────────────────────────────
\pagestyle{fancy}
\fancyhf{}
\renewcommand{\headrulewidth}{0.4pt}
\fancyhead[L]{\small\color{darkgray} Edge-Compute Pacing \& Feedback Node}
\fancyhead[R]{\small\color{darkgray} System Design Document}
\fancyfoot[L]{\small\color{darkgray} CPSC 498 / ECE Senior Project}
\fancyfoot[R]{\small\color{darkgray}\thepage}

% ── List style ─────────────────────────────────────────────────────────────────
\setlist[itemize]{noitemsep, topsep=2pt, leftmargin=1.4em}
\setlist[enumerate]{noitemsep, topsep=2pt, leftmargin=1.4em}

% ── Small helper for decision-record boxes ─────────────────────────────────────
\newcommand{\decision}[2]{%
  \par\vspace{0.4em}
  \noindent\colorbox{lightgray}{\parbox{0.97\linewidth}{\vspace{0.3em}
  \textbf{\color{trinityblue} Decision:} #1 \\[0.3em]
  \textbf{Rationale:} #2
  \vspace{0.3em}}}%
  \par\vspace{0.4em}
}

% ==============================================================================
\begin{document}

% ── Title Block ───────────────────────────────────────────────────────────────
\begin{center}
  {\LARGE\bfseries\color{trinityblue} Edge-Compute Pacing \& Feedback Node}\\[0.4em]
  {\large\color{darkgray} System \& Firmware Design Document}\\[0.8em]
  \begin{tabular}{rl}
    \textbf{Students:}          & Sean Balbale, Zoe Davidow \& Kayla Walsh \\
    \textbf{Program:}           & B.S.\ ECE\,/\,B.S.\ CS, Trinity College (Class of 2027) \\
    \textbf{Faculty Advisor:}   & Prof.\ Kenneth A.\ Kousen \\
    \textbf{Concept Summary:}   & April 28, 2026 \\
    \textbf{This Revision:}     & \parbox{3.6in}{September 2026 --- adds hardware/ECE design
                                   detail and revises the sensing model from power- to HR-based}
  \end{tabular}
\end{center}

\vspace{0.3em}
\noindent\rule{\linewidth}{1.2pt}
\vspace{0.3em}

% ── Abstract ──────────────────────────────────────────────────────────────────
\section{Abstract}

This project develops an embedded edge-compute node that gives endurance athletes ---
specifically rowers and cyclists --- real-time fatigue feedback with no phone and no cloud
in the loop. The device pairs with a wireless heart-rate chest strap (BLE, with an ANT+
fallback path if required), runs a physiologically motivated depletion/recovery model directly
on a resource-constrained microcontroller, and expresses the result through an on-board LED
array and buzzer. The contribution spans two disciplines: an \textbf{ECE} contribution in
sensor interfacing, power budgeting, and concurrent real-time firmware on a strict memory and
power envelope, and a \textbf{CS} contribution in the discretization and validation of a
continuous-time physiological model running on fixed-point, resource-constrained hardware.

% ── Problem Statement ─────────────────────────────────────────────────────────
\section{Problem Statement}

Commercial cycling and rowing head units display raw telemetry (watts, cadence, heart rate,
split) but perform no predictive computation. Athletes must self-interpret numbers while under
physiological stress, a cognitively expensive task that degrades pacing decisions. Smartphone
and cloud offload solve the compute problem but introduce connectivity dependency, latency, and
battery vulnerability --- unacceptable in competitive race settings.

\begin{center}
  \colorbox{lightgray}{\parbox{0.88\linewidth}{
      \vspace{0.4em}
      \textit{Can a real-time fatigue model be implemented with sub-100\,ms latency on a
        bare-metal or RTOS microcontroller, with sufficient numerical precision to provide
      actionable pacing feedback --- entirely offline, and driven by sensor hardware that is
      realistically available to a rower or cyclist without a \$400+ power meter?}
      \vspace{0.4em}
  }}
\end{center}

% ── Sensing and Fatigue Model ─────────────────────────────────────────────────
\section{Sensing \& Fatigue Model}

\subsection{From Power-Based to HR-Based Sensing}

The original concept used the $W'$ balance model, driven by an athlete's instantaneous power
output. In practice, the vast majority of rowers and club/collegiate cyclists do not own a
power meter (\$400--\$1500), while heart-rate chest straps are inexpensive (\$50--\$100) and
already owned by most endurance athletes. To make the device usable by an actual team rather
than only athletes with power-meter-equipped bikes, the primary sensor input is changed to
heart rate.

\decision{Primary fatigue input changes from power (W) to heart rate (bpm).}
{Sensor availability and cost across the target user base (teammates, club athletes) rather
than sensor precision. This is an access/equity decision, not a claim that HR is a better
signal than power.}

This has a real algorithmic consequence, documented for the advisor and for the written
thesis: the Monod--Scherrer/Skiba $W'$ model is defined and validated in terms of power against
Critical Power. Heart rate is not power. It lags true effort by tens of seconds to minutes, does
not register short anaerobic surges the way power does, and drifts with heat, hydration, and
accumulated fatigue independent of instantaneous output. Sections \ref{sec:wprime} and
\ref{sec:hprime} below keep the validated model as the literature reference case and present the
HR-driven model as an explicitly \textit{adapted, unvalidated} analog --- which becomes a
research question in its own right (Section \ref{sec:cs-questions}) rather than a
weakened version of the original claim.

\subsection{Reference Case: The $W'$ Balance Model}
\label{sec:wprime}

Two athlete-specific parameters govern the reference model:
\begin{itemize}
  \item $CP$ --- Critical Power (W), the highest sustainable steady-state output.
  \item $W'$ --- the finite anaerobic energy reservoir (J) above $CP$.
\end{itemize}

\begin{equation}
  W'_{\mathrm{bal}}(t) = W' - \int_{0}^{t} \bigl[P(\tau) - CP\bigr]^{+} \, d\tau
  \label{eq:depletion}
\end{equation}

\begin{equation}
  \frac{dW'_{\mathrm{bal}}}{dt} =
  \begin{cases}
    -\!\bigl(P(t) - CP\bigr) & P(t) \geq CP \quad (\text{depletion}) \\[6pt]
    \dfrac{W' - W'_{\mathrm{bal}}(t)}{\tau_{W'}} & P(t) < CP \quad (\text{recovery})
  \end{cases}
  \label{eq:skiba}
\end{equation}

\begin{equation}
  \tau_{W'}(t) = 546\,e^{-0.01\,[CP - P(t)]} + 316 \quad [\text{s}]
  \label{eq:tau}
\end{equation}

This remains the literature-grounded case used to validate the discretization approach
(Section \ref{sec:discrete}) whenever a power meter is available for bench testing.

\subsection{Adapted Case: HR-Analog Depletion Model}
\label{sec:hprime}

The primary, deployed model substitutes heart rate for power and lactate-threshold heart rate
($LTHR$) for Critical Power, keeping the same depletion/recovery structure:

\begin{equation}
  H'_{\mathrm{bal}}(t) = H' - \int_{0}^{t} \bigl[\mathrm{HR}(\tau) - LTHR\bigr]^{+} \, d\tau
  \label{eq:hdepletion}
\end{equation}

\begin{equation}
  \frac{dH'_{\mathrm{bal}}}{dt} =
  \begin{cases}
    -\!\bigl(\mathrm{HR}(t) - LTHR\bigr) & \mathrm{HR}(t) \geq LTHR \\[6pt]
    \dfrac{H' - H'_{\mathrm{bal}}(t)}{\tau_{H'}} & \mathrm{HR}(t) < LTHR
  \end{cases}
  \label{eq:hskiba}
\end{equation}

where $H'$ is a normalized reservoir size (not a physical energy quantity --- there is no
literature basis for an anaerobic joule-equivalent in HR terms) and $\tau_{H'}$ reuses the
functional form of Eq.~\eqref{eq:tau} with $LTHR$ in place of $CP$, pending empirical
re-fitting in Phase 2. This substitution is stated plainly as a design assumption to be tested,
not asserted as physiologically validated.

\subsection{Threshold Calibration ($LTHR$, $HR_{max}$, $HR_{rest}$)}

Three athlete-specific constants are required and must be enterable in the field (Section
\ref{sec:ui}):
\begin{itemize}
  \item \textbf{$LTHR$} --- primary threshold for Eq.~\eqref{eq:hskiba}. Estimated via a guided
    on-device field test (last 20 minutes of a 30-minute time-trial effort, averaged), with
    manual entry as an override for athletes with a known lab-tested value.
  \item \textbf{$HR_{max}$, $HR_{rest}$} --- used for Karvonen \%HRR normalization in the
    calibration UI and for sanity-bounding sensor readings (reject implausible HR samples).
\end{itemize}

\subsection{Discrete-Time Implementation}
\label{sec:discrete}

Both models share one parameterized Euler update, so the same firmware path handles the power
case (bench validation) and the HR case (deployed):

\begin{equation}
  X'_{\mathrm{bal}}[n] =
  \begin{cases}
    X'_{\mathrm{bal}}[n-1] - \bigl(S[n] - T\bigr)\,\Delta t & S[n] \geq T \\[8pt]
    X'_{\mathrm{bal}}[n-1] + \dfrac{X' - X'_{\mathrm{bal}}[n-1]}{\tau_{X'}[n]}\,\Delta t & S[n] < T
  \end{cases}
  \label{eq:discrete}
\end{equation}

where $(S, T, X') = (P, CP, W')$ for the reference case or $(\mathrm{HR}, LTHR, H')$ for the
deployed case, $\Delta t = 1\,\mathrm{s}$ (matching the BLE HR notification interval), and
$X'_{\mathrm{bal}}[0] = X'$. Fixed-point arithmetic only; well within the Cortex-M4F budget at
\SI{64}{\mega\hertz}.

\subsection{LED Feedback Mapping}

\begin{equation}
  k_{\mathrm{lit}}[n] = \left\lfloor N \cdot \frac{X'_{\mathrm{bal}}[n]}{X'} \right\rfloor,
  \qquad
  \mathrm{hue}[n] =
  \begin{cases}
    \text{green} & k_{\mathrm{lit}}/N > 0.50 \\
    \text{amber} & 0.20 < k_{\mathrm{lit}}/N \leq 0.50 \\
    \text{red}   & k_{\mathrm{lit}}/N \leq 0.20
  \end{cases}
\end{equation}

A buzzer interrupt fires when $X'_{\mathrm{bal}}[n] < 0.10 \cdot X'$.

% ── Hardware Architecture ─────────────────────────────────────────────────────
\section{Hardware Architecture \& Component Selection}
\label{sec:hardware}

\subsection{Microcontroller}

\decision{Nordic nRF52840 (Adafruit Feather nRF52840 Express for prototyping).}
{Native dual-radio silicon (BLE 5 and ANT, on the same die) means the sensor-protocol decision
below is a \emph{firmware} choice, not a board respin. The Feather's built-in USB and LiPo
charge circuitry also removes a discrete charger IC from the BOM. Cortex-M4F @ \SI{64}{\mega\hertz},
\SI{256}{\kilo\byte} RAM, \SI{1}{\mega\byte} flash --- ample headroom for the Euler-update
compute load, an I2C OLED driver, and FreeRTOS.}

\subsection{Wireless Sensor Protocol: BLE vs.\ ANT+}
\label{sec:protocol}

\begin{center}
\small
\begin{tabular}{>{\raggedright\arraybackslash}p{0.20\linewidth}>{\raggedright\arraybackslash}p{0.30\linewidth}>{\raggedright\arraybackslash}p{0.30\linewidth}}
\toprule
 & \textbf{BLE-only (Arduino/CircuitPython)} & \textbf{ANT+ (S340 SoftDevice)} \\
\midrule
Toolchain & Stays on Adafruit's Arduino core, current firmware plan unchanged & Requires leaving
Arduino for the nRF5 SDK / nRF Connect SDK; free ANT+ Adopter registration required to obtain
the S340 hex \\
Strap compatibility & Covers any dual-mode strap (Polar H10, Garmin HRM-Dual/Pro, Wahoo TICKR
--- confirmed to broadcast BLE + ANT+ simultaneously) & Required only for ANT+-only hardware
(rare among current chest straps) \\
Multi-athlete environments & Standard HR Service scan, connect to strap's stored MAC address & 
15 configurable channels, better suited to crowded-field broadcast scenarios (e.g.\ a regatta
with dozens of straps active) \\
\bottomrule
\end{tabular}
\end{center}

\decision{Default to BLE-only; keep ANT+ as a documented fallback path, not a Phase 1 requirement.}
{The straps realistically available to the team broadcast both protocols simultaneously, so
BLE-only satisfies Phase 1 without the SDK migration cost. The ANT+ path (Section
\ref{sec:protocol}, right column) is retained as a designed-for option if a strap without BLE
turns up, or if channel behavior in a crowded regatta start line makes ANT+ preferable ---
this is an open item, see Section \ref{sec:risks}.}

\subsection{Power System}
\label{sec:power}

\decision{3.7\,V LiPo pouch cell (JST-PH, \SI{2000}{\milli\ampere\hour}) instead of an 18650 cell
with a separate BMS/charger module.}
{The Feather's onboard charge IC and battery-voltage monitoring already expect exactly this
battery type and connector --- adding a bare 18650 would mean sourcing and wiring a separate
TP4056-class charger and physical cell holder for no functional benefit at this power budget.}

\begin{center}
\small
\begin{longtable}{p{0.30\linewidth}p{0.14\linewidth}p{0.10\linewidth}p{0.36\linewidth}}
\toprule
\textbf{Subsystem} & \textbf{Est.\ avg.\ current} & \textbf{Duty} & \textbf{Notes} \\
\midrule
nRF52840 core (sensor sampling + Euler compute) & \SI{5}{\milli\ampere} & continuous & Cortex-M4F
active during 1\,Hz compute tick, sleeps between ticks \\
BLE connection (HR notifications @ 1\,Hz) & \SI{1.5}{\milli\ampere} & continuous & Connection
interval tuned for the 1\,s sensor cadence, not BLE's fastest possible rate \\
NeoPixel strip (8 px, typical amber/green states) & \SI{15}{\milli\ampere} & continuous &
Worst case (8 px full white) is \textasciitilde\SI{144}{\milli\ampere}; normal operation stays
well under this \\
OLED (calibration menu only) & \SI{15}{\milli\ampere} & intermittent & Off during a session;
averages under \SI{1}{\milli\ampere} over a full ride/row \\
Buzzer (threshold alerts) & \SI{30}{\milli\ampere} & intermittent & Averages under
\SI{0.5}{\milli\ampere} over a session \\
\midrule
\textbf{Estimated average draw} & \textbf{\textasciitilde\SI{23}{\milli\ampere}} & & \\
\textbf{Estimated runtime on \SI{2000}{\milli\ampere\hour}} & \textbf{\textasciitilde 85\,h} & &
Far exceeds the \SI{6}{\hour} target \\
\bottomrule
\end{longtable}
\end{center}

\textit{These are first-pass estimates from component datasheets, not bench measurements ---
Phase 1 hardware bring-up includes a runtime validation pass under real load (Section
\ref{sec:timeline}).} Given the wide margin above the \SI{6}{\hour} target, a smaller
\SI{500}{\milli\ampere\hour}--\SI{1200}{\milli\ampere\hour} cell is worth evaluating once real
draw is measured, trading unused runtime for a smaller, lighter handlebar/hull-mount enclosure.

\subsection{User Interface Hardware}
\label{sec:ui}

\begin{itemize}
  \item \textbf{Output:} WS2812B NeoPixel stick (8 px) for the primary $X'_{\mathrm{bal}}$
    readout; piezo buzzer for the low-reserve alert.
  \item \textbf{Calibration input:} three 6\,mm tactile buttons (Up / Down / Select) plus a
    0.91-inch 128$\times$32 I2C OLED (STEMMA QT), added specifically because entering a
    3-digit $LTHR$ value via LED blink-count feedback is error-prone in the field. The OLED
    shares the I2C bus with no other peripheral, so it adds two GPIO pins total.
  \item \textbf{Config mode:} Select cycles $LTHR \rightarrow HR_{max} \rightarrow HR_{rest}
    \rightarrow$ save/exit; Up/Down adjust the active field with accelerating auto-repeat on a
    held press. Buttons are inert during an active session to prevent accidental input
    mid-effort.
\end{itemize}

\subsection{USB Connector}
\label{sec:usbc}

\decision{Add a USB Type-C breakout board (downstream configuration) wired to the Feather's
onboard Micro-USB connector, rather than using the stock Micro-USB port directly or an external
Micro-B-to-C adapter cable.}
{The Feather's native USB already handles both charging and full data (serial console, DFU
flashing) --- the only gap is that its physical connector is Micro-USB. An adapter cable would
close that gap with zero rework, but the breakout gives a cleaner result: it has the 5.1\,k$\Omega$
CC pulldown resistors needed for reliable 5\,V/1.5\,A negotiation built in (avoiding the
charger-compatibility question an unpowered adapter cable carries), and it panel-mounts flush
into the 3D-printed enclosure wall instead of leaving an internal cable loop.}

The nRF52840's USB D+/D$-$ lines are not broken out to any 0.1" header pin on the Feather ---
they run only to the pads of the onboard Micro-USB connector. Implementation is therefore a
piggyback solder job, not a header-wire connection: four fine-gauge (30\,AWG) wires are tacked
directly onto the Micro-USB connector's VBUS, D+, D$-$, and GND legs (the 5th, ID, pin is
unused) and routed to the matching pads on the breakout board. Desoldering the Micro-USB
connector outright is deliberately avoided, since lifting a pad on a fine-pitch SMT part is a
real risk without hot-air rework equipment. Continuity should be bench-verified with a
multimeter before final enclosure assembly, since the joints are not inspectable once sealed.

\subsection{Mechanical / Form Factor}

Handlebar mount (cycling) or hull-mount (rowing shell), 3D-printed in-house. Rowing introduces a
splash/spray exposure the original cycling-only form factor did not need to fully account for;
enclosure sealing around the button, OLED, and USB-C cutouts is called out as an open item
(Section \ref{sec:risks}).

% ── Firmware ───────────────────────────────────────────────────────────────
\section{Firmware \& Software Architecture}

\subsection{RTOS Task Structure}

FreeRTOS with four tasks on defined priorities and circular-buffer IPC: \textbf{sensor} (BLE HR
notification handling), \textbf{compute} (Eq.~\eqref{eq:discrete} Euler tick), \textbf{output}
(NeoPixel/buzzer drive), and \textbf{UI} (button debounce and OLED refresh during calibration).
The UI task is lowest priority and never blocks the sensor-to-LED path, preserving the
\SI{100}{\milli\second} end-to-end latency target from the original problem statement.

\subsection{Sensor Pairing}

\begin{itemize}
  \item \textbf{BLE:} scan for the standard Heart Rate Service (UUID 0x180D); connect to the
    first match; store the strap's MAC address in flash. Subsequent boots connect directly to
    the stored address rather than re-scanning, falling back to a fresh scan only on a
    connection timeout. This avoids picking up a nearby teammate's strap during team practice
    or a race start.
  \item \textbf{ANT+ (if implemented per Section \ref{sec:protocol}):} wildcard search on HR
    profile (device type 120) once, then lock to the discovered device number for all
    subsequent searches.
  \item Both paths expose a manual re-pair trigger (button hold) for swapping straps.
\end{itemize}

\subsection{Language \& Validation}

C/C++ against the nRF5 SDK (or Adafruit's Arduino core for the BLE-only path). The Euler-update
compute path is unit-tested against a reference Python implementation of Eq.~\eqref{eq:discrete}
before on-device deployment, for both the power and HR parameterizations.

% ── Research Framing ──────────────────────────────────────────────────────────
\section{Research Framing}

\subsection{Computer Science Research Questions}
\label{sec:cs-questions}

\begin{enumerate}
  \item \textbf{Algorithmic correctness under discretization.} How does Euler-step sampling
    frequency affect $X'_{\mathrm{bal}}$ accuracy relative to the continuous-time ground truth,
    for \emph{both} the validated power case and the adapted HR case?
  \item \textbf{Model transfer.} Does the depletion/recovery structure validated for power
    (Skiba et al.) produce a comparably informative fatigue signal when driven by heart rate
    instead --- and if the two disagree, where and why? (Requires bench sessions with both a
    power meter and HR strap on the same athlete for direct comparison.)
  \item \textbf{Concurrent firmware architecture.} What FreeRTOS scheduling policy minimizes
    worst-case latency while preventing sensor-buffer overflow on a single-core MCU?
  \item \textbf{Extensibility toward ML (stretch goal).} Can a lightweight recurrent model
    trained on historical HR traces forecast near-term HR well enough to make the LED feedback
    anticipatory rather than reactive, addressing the cardiac-lag weakness named in Section
    \ref{sec:hprime}? Full specification in Section \ref{sec:ml-stretch}.
\end{enumerate}

\subsection{ECE Research Questions}

\begin{enumerate}
  \item \textbf{Power budget validation.} Does measured current draw under real BLE connection
    and NeoPixel activity match the datasheet-derived estimate in Section \ref{sec:power}, and does the
    system meet the \SI{6}{\hour} runtime target with margin to downsize the battery?
  \item \textbf{RF behavior in the deployed environment.} Does BLE connection stability hold up
    mounted on a carbon bicycle frame or an aluminum/carbon rowing shell (both RF-unfriendly
    environments), and does antenna placement in the 3D-printed enclosure need adjustment?
  \item \textbf{End-to-end latency budget.} With the wireless link now in the loop (strap
    $\rightarrow$ BLE notification $\rightarrow$ compute $\rightarrow$ LED), where does the
    \SI{100}{\milli\second} budget actually get spent, and is BLE connection-interval tuning
    sufficient to hold it?
\end{enumerate}

% ── ML Stretch Goal ────────────────────────────────────────────────────────────
\section{Stretch Goal: On-Device HR Forecasting}
\label{sec:ml-stretch}

This is scoped as a \textbf{Phase 5 stretch goal}, not a Phase 1--4 dependency --- the core
deliverable (Sections 3--6) is complete and demoable without it. It is included here because it
turns the model's stated weakness into the project's most interesting extension, and because a
real training dataset already exists (Section \ref{sec:ml-data}).

\subsection{Motivation}

Section \ref{sec:hprime} names cardiac lag as the central limitation of the HR-analog model: by
the time $\mathrm{HR}(t)$ crosses $LTHR$, the athlete's actual effort crossed it tens of seconds
earlier. The LED feedback is therefore always a step behind reality. A small forecasting model
sitting in front of the depletion equation --- predicting $\mathrm{HR}(t + \delta)$ from the
recent trace rather than reacting to $\mathrm{HR}(t)$ --- gives the athlete a genuine early
warning instead of a delayed one, which is a materially different (and more useful) product
behavior, not just a numerical refinement.

\subsection{Model \& Integration}

\begin{itemize}
  \item \textbf{Architecture:} single-layer GRU or LSTM, 8--16 hidden units, taking the trailing
    30--60 samples of $\mathrm{HR}(t)$ (plus first difference) at the existing 1\,Hz cadence and
    predicting $\mathrm{HR}(t + 5\text{--}10\,\mathrm{s})$.
  \item \textbf{Footprint:} at this size, a handful of hundred to $\sim$2000 parameters ---
    low single-digit \si{\kilo\byte} of flash at 8-bit quantization. RAM is dominated by the
    inference runtime overhead (TFLite Micro's interpreter, roughly
    \SIrange{2}{10}{\kilo\byte} depending on configuration) rather than the model itself, and
    should be benchmarked against whatever the S140/S340 SoftDevice reserves before assuming
    it fits (Section \ref{sec:power}'s RAM headroom is for the compute task, not the radio
    stack).
  \item \textbf{Inference cost:} a model this small runs in microseconds to low single-digit
    milliseconds on the Cortex-M4F --- does not threaten the \SI{100}{\milli\second} end-to-end
    latency budget (Section \ref{sec:discrete}).
  \item \textbf{Integration point:} the compute task (Section \ref{sec:discrete}) gains an
    inference substep ahead of the Euler tick; $H'_{\mathrm{bal}}$ is computed against the
    forecast $\widehat{\mathrm{HR}}(t+\delta)$ rather than the raw sample, and the LED hue
    transition fires on the forecast crossing $LTHR$.
\end{itemize}

\subsection{Training Data}
\label{sec:ml-data}

Unlike a typical student ML project, this does not need a bootstrapped or public dataset ---
historical session data already exists from AstrapheAI, Sean's self-hosted endurance-coaching
platform, which ingests Strava/WHOOP telemetry per workout including \textbf{power, heart rate,
speed, elevation}, and additional derived fields. This is a real, already-labeled corpus of
multi-athlete session data rather than a synthetic or public substitute, and is queryable
directly (workout listing and windowed stream extraction) rather than requiring a new export
pipeline.

\begin{itemize}
  \item \textbf{Session split:} train/validation/test split by athlete and by session (not by
    sample) to avoid leaking a single ride's autocorrelated HR trace across splits.
  \item \textbf{Label quality check before committing to this scope:} confirm sampling rate and
    gap frequency in the existing HR streams are consistent enough for a 1\,Hz forecasting
    target, and confirm whether per-session effort/RPE context exists or whether HR and power
    alone are the only signal available for training.
  \item \textbf{Cross-signal opportunity:} since power is also present in the same dataset for
    sessions where it was recorded, the same corpus supports the model-transfer research
    question (Section \ref{sec:cs-questions}, item 2) --- training and evaluating the forecast
    model against sessions where \emph{both} power and HR exist gives a direct check of how much
    the HR-analog model's lag actually costs, using the athlete's own data rather than a
    literature estimate.
  \item \textbf{Offline training:} standard Python/PyTorch or TensorFlow training loop, exported
    and quantized to TFLite Micro (or a small hand-rolled fixed-point inference routine, given
    the model's size) for on-device deployment; unit-tested against the Python reference
    forecast output before flashing, following the same validation pattern already used for the
    Euler update (Section \ref{sec:discrete}).
\end{itemize}

\subsection{Risk}

Genuinely a stretch goal: if Phase 1--4 timeline slips, this is the first thing to cut without
weakening the core deliverable, and the fallback in CS Research Question 4's original framing
(online $\tau_{H'}$ personalization via a lightweight estimator rather than a trained network)
remains available as a lower-effort substitute.

% ── Deliverables & Timeline ───────────────────────────────────────────────────
\section{Deliverables \& Timeline}
\label{sec:timeline}

\begin{center}
\begin{longtable}{cp{0.68\linewidth}l}
\toprule
\textbf{Phase} & \textbf{Milestone} & \textbf{Target} \\
\midrule
1 & Hardware bring-up; BLE HR strap pairing verified; power budget bench-validated & Sep 2026 \\
2 & $H'$ model on MCU; validated against Python reference and, where possible, against the
    power-based $W'$ reference case & Oct 2026 \\
3 & FreeRTOS integration; UI/calibration task; latency and scheduling benchmarks & Nov 2026 \\
4 & Full feedback loop; field test with HR strap on-bike and on-water & Dec 2026 \\
5 & Formal write-up, poster, and final public demo & Apr 2027 \\
5$^\dagger$ & \textit{Stretch:} HR-forecasting model (Section \ref{sec:ml-stretch}) --- training data
    audit, offline training, on-device integration & Apr 2027 \\
\bottomrule
\end{longtable}
{\small $^\dagger$Stretch goal; does not block the core Phase 5 deliverable if timeline is tight.}
\end{center}

% ── Bill of Materials ──────────────────────────────────────────────────────────
\section{Bill of Materials}
\label{sec:bom}

Linked items are current as of this revision; prices are list price at the linked retailer and
will drift. Items marked \textit{(source locally)} are generic enough that any supplier works
and are not worth pinning to a specific link.

\begin{center}
\small
\begin{longtable}{p{0.38\linewidth}p{0.10\linewidth}p{0.42\linewidth}}
\toprule
\textbf{Item} & \textbf{Price} & \textbf{Notes} \\
\midrule
\endhead
\href{https://www.adafruit.com/product/4062}{Adafruit Feather nRF52840 Express} &
  \$24.95 & Core MCU. Native BLE, USB, onboard LiPo charge circuit. \\
\href{https://www.adafruit.com/product/2011}{Lithium Ion Battery --- 3.7\,V \SI{2000}{\milli\ampere\hour}, JST-PH} &
  \$12.50 & Plugs directly into the Feather's onboard battery port; no separate charger/BMS needed. \\
\href{https://www.adafruit.com/product/1426}{NeoPixel Stick --- 8 x WS2812/SK6812 RGB LED} &
  \$5.95 & Primary $H'_{\mathrm{bal}}$ readout. \\
\href{https://www.adafruit.com/product/1536}{Buzzer 5V --- Breadboard friendly} &
  \$1.95 & Low-reserve alert; self-driving, apply 3--5\,V directly. \\
\href{https://www.adafruit.com/product/367}{Tactile Button Switch (6mm) --- 20 Pack} &
  \$4.95 & Up / Down / Select calibration buttons (3 used, 17 spares). \\
\href{https://www.adafruit.com/product/4440}{Monochrome 0.91" 128$\times$32 I2C OLED --- STEMMA QT} &
  $\sim$\$15 & Threshold-entry and live-HR display during calibration. \\
\href{https://www.adafruit.com/product/4209}{STEMMA QT to Premium Male Headers Cable} &
  \$0.95 & Solderless I2C hookup from OLED to Feather header pins. \\
\href{https://www.adafruit.com/product/4090}{USB Type C Breakout Board --- Downstream Connection} &
  \$2.95 & External USB-C port (Section \ref{sec:usbc}); wire to the Feather's Micro-USB pins. \\
\midrule
Protoboard / perfboard \textit{(source locally)} & $\sim$\$5 & \\
22--26\,AWG hookup wire, incl.\ fine 30\,AWG for the USB-C piggyback joints \textit{(source locally)} &
  $\sim$\$8 & \\
JST-PH 2-pin connectors, a few \textit{(source locally)} & $\sim$\$5 & For detachable buzzer/NeoPixel leads. \\
\midrule
\textbf{Linked components subtotal} & \textbf{$\sim$\$69} & \\
\bottomrule
\end{longtable}
\end{center}

\textbf{Optional --- HR chest strap for bench testing:}
\href{https://www.polar.com/us-en/sensors/h10-heart-rate-sensor}{Polar H10 Heart Rate Sensor}
($\sim$\$90--105) if the team does not already have a spare dual-mode (BLE + ANT+) strap
available for dedicated bench work. Any teammate's Polar H10, Garmin HRM-Dual/Pro, or Wahoo
TICKR works as-is with the BLE-only plan in Section \ref{sec:protocol}.

% ── Open Decisions & Risks ────────────────────────────────────────────────────
\section{Open Decisions \& Risks}
\label{sec:risks}

\begin{itemize}
  \item \textbf{ANT+ vs.\ BLE-only:} pending confirmation that every strap the team will
    actually use is dual-mode. If not, the ANT+/S340 path in Section \ref{sec:protocol} needs
    to move from "documented fallback" to "Phase 1 requirement."
  \item \textbf{$LTHR$ field-test accuracy:} the guided on-device test is a standard field
    estimate, not a lab-grade measurement; needs validation against a known value for at least
    one athlete before being trusted as the default calibration path.
  \item \textbf{Enclosure sealing:} rowing's splash exposure is a materially different
    environment than cycling for the button/OLED cutouts; needs a sealing approach before the
    on-water field test in Phase 4.
\end{itemize}

% ── Advisor Engagement ────────────────────────────────────────────────────────
\section{Advisor Engagement}

Prof.\ Kousen's advising continues to focus on software architecture, algorithmic design, and
CS rigor per the original concept summary; the hardware/ECE decisions in Section \ref{sec:hardware} are outside
his stated area and are tracked here for the ECE half of the advising/grading structure. Monthly
check-ins continue to cover the written deliverables (proposal, interim report, thesis) and, per
Section \ref{sec:cs-questions}, the model-transfer question this revision introduces.

\vspace{1em}
\noindent\rule{\linewidth}{0.5pt}
\vspace{0.2em}
{\small
  \noindent \textbf{Contact:}\;
  Sean Balbale ---
  \href{mailto:sean.balbale@trincoll.edu}{sean.balbale@trincoll.edu}
  $\;\big|\;$
  Zoe Davidow ---
  \href{mailto:zoe.davidow@trincoll.edu}{zoe.davidow@trincoll.edu}
  $\;\big|\;$
  Kayla Walsh ---
  \href{mailto:kayla.walsh@trincoll.edu}{kayla.walsh@trincoll.edu}
}

\end{document}